Proofs that a set of equations between circuits is \emph{complete} --- that it
derives every true equation --- are among the most computation-heavy proofs in
the theory of quantum computing: they hinge on normal forms, and on relentless
case analyses showing that every rule application interacts correctly with
every normal form.  Recent completeness proofs come with appendices hundreds
of pages long, produced and checked by hand.  This paper presents an Agda
library that mechanizes the mathematics these proofs rest on: monoid and group
presentations by generators and relations, normal forms, and completeness of
finite relation sets, specialized to wire-indexed circuits.  The library
represents presented monoids as inductively generated congruences on free
monoids of words, represents normal forms as the standard-library function
bundles (injections, right inverses, bijections) out of the word setoid, and
derives completeness from soundness and a \emph{unique normal form} witness by
a single generic lemma.  Its engine is a constructive, monoid-level rendition
of coset enumeration in the style of Reidemeister and Schreier: towers of
verified coset tables that assemble normal forms level by level, in the manner
of a mixed-radix number system, together with presentation-level constructions
--- direct, semidirect, and amalgamated free products, and group extensions ---
whose semantic counterparts we also contribute, targeting the Agda standard
library.  We validate the design on a catalogue of verified presentations:
symmetric groups as circuits (sound and complete for both an endofunction and
a permutation semantics), cyclic groups, wreath products $\mathbb{Z}_N \wr S_n$
(signed permutations), Pauli groups in all odd prime dimensions, and, as
amalgamated free products with verified alternating normal forms, the
single-qubit and single-qutrit Clifford+$T$ monoids and the matrix group
$U_3(\mathbb{Z}[\tfrac12,i])$.  All results are collected in a single index
module; the whole development is safe, axiom-free Agda.  We compare in detail
with formalizations of group presentations, normal forms, and quantum circuits
in Agda, Lean, Coq, and Isabelle.
